\section{Related Work}
\label{sec:related-work}

\subsection{Concept Erasure in Image and Video Diffusion}
\label{subsec:related-concept-erasure}

\outlineblock{Argument}{
Summarize training-free, fine-tuning, and representation-level erasure by the
unit they suppress. End with the precise gap: most treat the target as a
lexical or visual concept and do not distinguish the entity's role in a
physical event. Move method-by-method detail to
\appref{app:training-baselines}.}

\subsection{Video Object Removal and Controllable Editing}
\label{subsec:related-video-editing}

\outlineblock{Argument}{
Position object removal, effect-aware removal, and temporally consistent video
editing. Distinguish their input-video, mask, trajectory, or reference-frame
assumptions from open-ended text-to-video generation of a distributional
source-free world.}

\subsection{Counterfactual Supervision and Role-Based Generalization}
\label{subsec:related-counterfactual}

\outlineblock{Argument}{
Connect paired supervision, counterfactual targets, structured prompt
interventions, and role generalization. State explicitly that the paper tests
behavior induced by a registered source role; it does not claim to discover a
true internal causal representation.}

\claimboundary{Use a priority word such as \emph{first} only after the
literature search and citation ledger establish the exact scope of the claim.}
